\clearpage
\section{RX-side Conversion Gain --- Low-Side LO (No Spectral Flip)}
\label{sec:rx-gain-lowside}

\textbf{What is conversion gain?} The RX side is the TX side in reverse: it takes a signal in at a high FR3
frequency and brings it down to FR1, amplifying it along the way. Conversion gain measures that amplification
across the change of frequency. Feed a signal into the ``RX-side FR3 RF In'' connector at, say, −50\,dBm, and
measure the signal that comes out of the ``RX-side FR1 IF Out'' connector at its new frequency, say −5\,dBm. The conversion
gain is simply the difference: −5 − (−50) = 45\,dB. Because both powers are in dBm, subtracting them gives the
conversion gain in dB directly.

\textbf{What's plotted.} Each of the first seven figures holds FR1 (the IF output frequency) fixed at one
center value (1.0 up to the unit's 4.0\,GHz maximum, in 0.5\,GHz steps) and plots conversion gain against the
swept FR3 input frequency for that output, using the same methodology as the TX conversion-gain sections: every
point is the average conversion gain across a real 500\,MHz swept window centered on that point, not a single
CW tone.\footnote{The FieldFox's built-in tracking generator drives FR3 into the ``RX-side FR3 RF In'' port at a
calibrated
−50\,dBm while its own receiver sweeps the corresponding FR1 window at the ``RX-side FR1 IF Out'' port
in lock-step.} Grey is the raw LNA-and-mixer chain with the AM3153 filter bypassed; blue is the filter engaged,
re-selected at every point using the same ``optimal'' algorithm described in
Section~\ref{sec:sideband-suppression}, reused here because the AM3153 is RX's actual image-rejection mechanism
--- there is no I/Q image rejection on this side of the board, since the mixer's Q output is simply discarded.

\textbf{Measurement conditions.} \lmxDriveRule{} (set from the LO frequency). Because the LO is shared, the TX side was set up at every point exactly as in normal use (its stored LO-leakage calibration applied and its optimal filter engaged), with both TX-side ports terminated in 50\,\textohm. The unit was measured hot, as on the TX side: no fan, LMX2820 die at \condRxGainMinC--\condRxGainMaxC\textdegree C.

\textbf{How the filter is chosen.} The same two-step rule used throughout this datasheet: first find whichever
of the 48 candidate configurations (3 bands × 16 sub-states) passes the most desired energy on its own;
then, restricted to candidates within 3\,dB of that best desired power, choose whichever rejects the image the
most. Consecutive points that land on the same filter form a \emph{region}; the dashed vertical lines on each
panel mark every boundary between regions, and the blue hex code running up from the x-axis inside each region
is the actual \texttt{filter\_bits} value used for every point in it.

\textbf{Low-side LO: LO = FR3 − FR1.} This is the LO tuning regime shown in this section. Because of it, the
IF output spectrum is \emph{not} flipped relative to the RF input --- as FR3 increases, the IF output's own
internal content increases too.\footnote{Since TX and RX share one physical LO, this also means the two
chains' LO-side states are always consistent with each other: whenever RX is running low-side LO, TX (if active at the
same time) is necessarily running upper-sideband up-conversion, and vice versa.} In Figure~\ref{fig:7-7} (FR1 = 4.0\,GHz), as on the TX side, the window extends 250\,MHz past the unit's IF design spec. Figures~\ref{fig:7-8} and~\ref{fig:7-9} repeat the bypassed and optimal data across the full (FR1, FR3) grid.

\figpair{images/sec7_fig1.pdf}{FR1 (output) = 1.0\,GHz.}{fig:7-1}%
        {images/sec7_fig2.pdf}{FR1 (output) = 1.5\,GHz.}{fig:7-2}

\figpair{images/sec7_fig3.pdf}{FR1 (output) = 2.0\,GHz.}{fig:7-3}%
        {images/sec7_fig4.pdf}{FR1 (output) = 2.5\,GHz.}{fig:7-4}
\figpair{images/sec7_fig5.pdf}{FR1 (output) = 3.0\,GHz.}{fig:7-5}%
        {images/sec7_fig6.pdf}{FR1 (output) = 3.5\,GHz.}{fig:7-6}
\figonesmall{images/sec7_fig7.pdf}{FR1 (output) = 4.0\,GHz. 250\,MHz of this window falls outside the unit's IF design spec (shown in red).}{fig:7-7}

\newpage
\figrow{images/sec7_fig8.pdf}{Conversion gain, low-side LO: full measured (FR1, FR3) grid, filter bypassed.
Zoom in to read each cell's exact values, printed as (FR1, FR3, conversion gain), with FR1 and FR3 in GHz and
the conversion gain in dB.}{fig:7-8}
\figrow{images/sec7_fig9.pdf}{Conversion gain, low-side LO: full measured (FR1, FR3) grid, optimal filter.
Zoom in to read each cell's exact values, printed as (FR1, FR3, conversion gain), with FR1 and FR3 in GHz and
the conversion gain in dB.}{fig:7-9}
